% ============================================================
%  TFM Memoria: Pipeline Licitaciones
%  Estilo suizo: todo sans, monocromo, jerarquía tipográfica
%  Compila con XeLaTeX (fontspec). En Overleaf: Compiler → XeLaTeX
% ============================================================
\documentclass[11pt,a4paper]{article}

% ---------- Idioma ----------
\usepackage[spanish,es-tabla,es-noshorthands]{babel}
\usepackage{microtype}

% ---------- Tipografía: todo sans (XeLaTeX / fontspec) ----------
\usepackage{fontspec}
\IfFontExistsTF{Roboto}{
  \setmainfont{Roboto}
  \setsansfont{Roboto}
}{
  \setmainfont{Noto Sans}
  \setsansfont{Noto Sans}
}
\IfFontExistsTF{Roboto Mono}{
  \setmonofont{Roboto Mono}[Scale=MatchLowercase]
}{
  \setmonofont{Noto Sans Mono}[Scale=MatchLowercase]
}
\renewcommand{\familydefault}{\sfdefault}

% ---------- Color: solo grises ----------
\usepackage[table]{xcolor}
\definecolor{ink}{HTML}{111111}    % texto
\definecolor{muted}{HTML}{6B6B6B}  % secundario
\definecolor{line}{HTML}{CFCFCF}   % reglas
\definecolor{paper2}{HTML}{F4F4F4} % fondo suave
\definecolor{enlace}{HTML}{1E40AF} % azul sobrio solo para enlaces

% ---------- Página ----------
\usepackage[a4paper,top=2.5cm,bottom=2.7cm,left=2.6cm,right=2.6cm,headsep=16pt]{geometry}
\usepackage{parskip}
\setlength{\parskip}{0.5em}

% ---------- Cabecera / pie ----------
\usepackage{fancyhdr}
\pagestyle{fancy}
\setlength{\headheight}{13pt}
\fancyhf{}
\renewcommand{\headrulewidth}{0.4pt}
\renewcommand{\headrule}{\color{line}\hrule width\headwidth height 0.4pt}
\fancyhead[L]{\footnotesize\color{muted}TFM · Pipeline de contratación pública}
\fancyhead[R]{\footnotesize\color{muted}\nouppercase{\leftmark}}
\fancyfoot[L]{\footnotesize\color{muted}T. Campoy Rojo}
\fancyfoot[R]{\footnotesize\bfseries\thepage}

% ---------- Secciones ----------
\usepackage{titlesec}
\titleformat{\section}
  {\bfseries\Large}{\thesection}{0.7em}{}
  [{\color{line}\vspace{3pt}\titlerule[0.6pt]}]
\titlespacing*{\section}{0pt}{2em}{1em}
\titleformat{\subsection}
  {\bfseries\large}{\thesubsection}{0.6em}{}
\titlespacing*{\subsection}{0pt}{1.4em}{0.5em}
\titleformat{\subsubsection}[runin]
  {\bfseries\color{muted}}{\thesubsubsection}{0.5em}{}

% ---------- Tablas ----------
\usepackage{booktabs}
\usepackage{tabularx}
\usepackage{array}
\renewcommand{\arraystretch}{1.3}
\newcommand{\cabecera}[1]{{\bfseries #1}}

% ---------- Figuras ----------
\usepackage{graphicx}
\usepackage{caption}
\captionsetup{font={small},labelfont={bf},labelsep=quad}
\usepackage{float}
\usepackage{placeins}
\usepackage{needspace}
\usepackage{tikz}
\usetikzlibrary{arrows.meta,positioning,calc,fit,backgrounds}

% ---------- Código ----------
\usepackage{listings}
\lstdefinestyle{plain}{
  basicstyle=\ttfamily\footnotesize\color{ink},
  backgroundcolor=\color{paper2},
  keywordstyle=\bfseries,
  commentstyle=\color{muted}\itshape,
  stringstyle=\color{ink},
  showstringspaces=false,
  columns=fullflexible,
  frame=none, xleftmargin=12pt, xrightmargin=8pt,
  aboveskip=1em, belowskip=1em
}
\lstset{style=plain}

% ---------- Listas ----------
\usepackage{enumitem}
\setlist[itemize]{leftmargin=1.3em, itemsep=0.3em, topsep=0.45em,
  label=\textcolor{muted}{\rule[0.35ex]{4pt}{4pt}}}
\setlist[enumerate]{leftmargin=1.6em, itemsep=0.3em, topsep=0.45em}

% ---------- Enlaces ----------
\usepackage{hyperref}
\hypersetup{colorlinks=true, linkcolor=enlace, urlcolor=enlace, citecolor=enlace,
  pdftitle={TFM: Pipeline reproducible para datos de contratación pública},
  pdfauthor={Tomás Campoy Rojo}}

% ---------- Utilidades ----------
\newcommand{\tarrow}{$\rightarrow$}

\setcounter{secnumdepth}{2}
\setcounter{tocdepth}{2}
\renewcommand{\contentsname}{Contenido}
\addto\captionsspanish{\renewcommand{\refname}{Bibliografía}}

\begin{document}

% ============================================================
%  PORTADA
% ============================================================
\begin{titlepage}
\raggedright
{\small\bfseries\MakeUppercase{Trabajo Fin de Máster}}\\[3pt]
{\small\color{muted}Máster en Ingeniería de Datos · Universidad Complutense de Madrid}\\[1em]
{\color{ink}\rule{\textwidth}{0.8pt}}

\vspace{4cm}

{\fontsize{32}{37}\selectfont\bfseries
Pipeline reproducible para integrar y analizar datos de contratación pública\par}

\vspace{1.2cm}

{\large\color{muted}TED · Open Data Portal de la Plaza (PLACSP) · DIR3 · CPV\par}

\vfill

{\color{line}\rule{\textwidth}{0.4pt}}\\[1em]
{\bfseries Tomás Campoy Rojo}\\[2pt]
{\color{muted}Autor}\\[0.9em]
{\bfseries Septiembre de 2026}
\end{titlepage}

% ============================================================
%  ÍNDICE
% ============================================================
\tableofcontents
\clearpage

% ============================================================
%  1. INTRODUCCIÓN
% ============================================================
\section{Introducción}
\markboth{Introducción}{Introducción}

Cada año, las administraciones públicas españolas y europeas publican decenas de
miles de oportunidades de negocio en portales oficiales heterogéneos. La
información existe, es abierta y gratuita, pero llega fragmentada en formatos
distintos (JSON, ZIP, Atom, CODICE/XML), con reglas de publicación y revisión
propias. Para una empresa que vende tecnología, la diferencia entre encontrar
una licitación a tiempo y descubrirla cuando ya ha finalizado el plazo es,
puede suponer perder una oportunidad de negocio.

Este TFM aborda ese problema desde la ingeniería de datos: construir un
pipeline reproducible que \textbf{integre, normalice y estructure} los avisos de
licitación publicados en TED (ámbito europeo) y en el portal de datos abiertos
de la Plataforma de Contratación (OpenPLACSP), dejando el dato listo para
análisis y consumo.

Integrar ambas fuentes no es un mero ejercicio de formato. Cada portal usa sus
propios identificadores, sus reglas de versionado y su granularidad temporal:
un expediente puede aparecer como aviso previo, rectificarse varias veces,
adjudicarse o anularse, y cada uno de esos hechos llega como una publicación
distinta que hay que reconciliar sin perder la traza.

También hay que registrar qué periodo cubre cada descarga, si está completo y
qué contenido exacto se obtuvo. Sin esa procedencia, dos ejecuciones no se
pueden comparar ni auditar.

Para una empresa proveedora, el valor está en pasar de portales fragmentados a
una vista accionable: qué procedimientos siguen abiertos, con qué plazo, en
qué categorías CPV encajan con su oferta y qué historial tiene el comprador.
El resultado del proyecto es un dashboard analítico que permite consultar esa
información sin conocer el funcionamiento interno del pipeline.

\subsection{Pregunta de investigación}

\emph{¿Hasta qué punto es posible construir con fuentes abiertas oficiales un
histórico unificado y trazable de la contratación pública, con identidad
estable de los avisos ante revisiones y anulaciones, que soporte productos de
datos como un feed diario de oportunidades?}

\subsection{Contribuciones}

\begin{itemize}
  \item Un pipeline por capas (raw \tarrow{} Bronze \tarrow{} Silver \tarrow{} Gold) con procedencia
        explícita y rechazos localizados en lugar de descartes silenciosos.
  \item Un modelo canónico de eventos de contratación que preserva el histórico
        de avisos, revisiones y \emph{tombstones} con identidad determinista.
  \item Un estado vigente derivado con Spark desde el histórico canónico, con
        política de apertura congelada sobre el ciclo de vida oficial de PLACSP
        (\texttt{open\_opportunities} en Parquet).
  \item Una capa analítica ligera (DuckDB) que expone vistas reproducibles y un
        export determinista a CSV para Tableau, sin redefinir la lógica de
        negocio.
  \item Un contrato de datos versionado entre capas, junto con un \emph{manifest}
        de ejecución con conteos y checksums.
  \item Un marco de evaluación reproducible (mismo resultado y
        rendimiento) como base para decisiones de implementación: las mejoras
        solo se aceptan si mantienen el comportamiento esperado.
\end{itemize}

% ============================================================
%  2. OBJETIVOS Y ALCANCE
% ============================================================
\section{Objetivos y alcance}

El alcance se divide en un núcleo obligatorio y varias extensiones, de modo
que el resultado principal sea entregable dentro del plazo del máster.

\subsection{Núcleo (entregable mínimo)}

\begin{itemize}
  \item Descarga explícita y conservación local de TED (API JSON), OpenPLACSP
        (ZIP, Atom, CODICE/XML) y foto DIR3 de unidades orgánicas.
  \item Parsing y normalización a Bronze Parquet con procedencia, rechazos
        localizados y conteos de ingesta.
  \item Silver canónico: \texttt{procurement\_events.parquet} con histórico de
        avisos, revisiones y \emph{tombstones} como eventos, construido con
        Polars y verificado contra la implementación inicial.
  \item Gold: estado vigente resuelto con Spark desde Silver canónico
        (\texttt{open\_opportunities} en Parquet), con política
        abierto/accionable congelada sobre el ciclo de vida oficial y
        enriquecimiento CPV/DIR3; la vista legada en memoria (TenderRecord) y el
        Gold provisional quedaron sustituidos.
  \item Capa analítica DuckDB (vistas versionadas solo-lectura sobre Gold) y
        export determinista a CSV para el dashboard Tableau.
  \item Contrato de datos versionado y \emph{manifest} de ejecución con checksums.
\end{itemize}

\subsection{Extensiones (si hay tiempo)}

\begin{itemize}
  \item Orquestación programada con Airflow sobre las funciones actuales del
        pipeline. El empaquetado reproducible con Docker ya está implementado.
  \item Feed diario de oportunidades con filtros por perfil de proveedor.
\end{itemize}

\subsection{Criterios de éxito}

El trabajo se considera completo cuando la cadena puede ejecutarse de forma
reproducible desde los datos de entrada hasta los conjuntos utilizados por
Tableau, los registros rechazados quedan identificados con su causa y el
histórico de revisiones puede reconstruirse desde Silver. La ejecución debe
dejar conteos, checksums de las salidas deterministas y artefactos de control
suficientes para comprobar el resultado.

% ============================================================
%  3. FUENTES DE DATOS
% ============================================================
\section{Fuentes de datos}

\subsection{Panorama}

\begin{table}[ht]
\centering
\small
\begin{tabularx}{\textwidth}{@{}l l l X@{}}
\toprule
\cabecera{Fuente} & \cabecera{Ámbito} & \cabecera{Formato} & \cabecera{Papel en el pipeline}\\
\midrule
TED eSender & Europeo & JSON (API) & Avisos previos y adjudicaciones de toda la UE\\
Open Data PLACSP & España & ZIP · Atom · CODICE/XML & Licitaciones y adjudicaciones nacionales\\
DIR3 & España & CSV/online & Dimensión de unidades orgánicas por ámbito\\
CPV & UE & Reglamento 2015/2424 & Taxonomía oficial de objeto del contrato\\
\bottomrule
\end{tabularx}
\caption{Fuentes oficiales integradas y su función en la plataforma.}
\end{table}

El corpus de referencia cubre enero--junio de 2026: seis ZIP mensuales de
OpenPLACSP y un JSONL de TED, con rutas, tamaños y checksums fijados en el
manifest versionado. La consulta TED de ese corpus incluye términos
tecnológicos, luego está sesgada al caso de estudio y no representa el
histórico completo.

La dimensión de compradores se apoya en DIR3, el directorio común de unidades
orgánicas: seis ámbitos con 108.261 unidades en la foto de septiembre de 2026.

\begin{table}[ht]
\centering
\small
\begin{tabularx}{0.62\textwidth}{@{}X r@{}}
\toprule
\cabecera{Ámbito DIR3} & \cabecera{Unidades}\\
\midrule
AGE & 19.919\\
CCAA & 18.825\\
EELL & 26.606\\
Universidades & 7.360\\
Otras instituciones & 1.338\\
Justicia & 34.213\\
\midrule
\cabecera{Total} & \cabecera{108.261}\\
\bottomrule
\end{tabularx}
\caption{DIR3 por ámbito (foto de septiembre de 2026).}
\end{table}

La taxonomía de objeto del contrato es el vocabulario CPV 2008: 9.454 códigos
en un cuerpo de 8 dígitos más un noveno dígito de control tras un guion
(\texttt{XXXXXXXX-Y}). El dígito de control solo detecta erratas y los ceros a
la izquierda son significativos: la clave canónica es el cuerpo como cadena,
nunca como entero.

\subsection{Volumen del corpus}

El corpus de referencia reúne datos de TED y OpenPLACSP entre enero y junio de
2026. Los seis ZIP mensuales de OpenPLACSP ocupan 1.118.920.855 bytes y el
fichero JSONL de TED, 105.503.327 bytes. Los tamaños y checksums individuales
se conservan en \path{data/gold/run_manifest.json}.

\begin{table}[ht]
\centering
\small
\begin{tabularx}{\textwidth}{@{}l X r r@{}}
\toprule
\cabecera{Fuente} & \cabecera{Periodo o ficheros} & \cabecera{Tamaño} & \cabecera{Entradas}\\
\midrule
OpenPLACSP & Enero--junio de 2026, seis ZIP & 1.118.920.855 bytes & 349.399\\
TED & Enero--junio de 2026, un JSONL & 105.503.327 bytes & 9.905\\
DIR3 & Foto de septiembre de 2026, seis ámbitos & n/d & 108.261\\
CPV 2008 & Dimensión oficial versionada & n/d & 9.454\\
\bottomrule
\end{tabularx}
\caption{Volumen de las fuentes y dimensiones del corpus de referencia.}
\end{table}

Las fuentes difieren tanto en su empaquetado como en su granularidad. El
pipeline conserva esos originales y los transforma en representaciones
Parquet comunes para que las etapas posteriores no dependan del formato de
cada portal. Los conteos de las capas derivadas se presentan en Resultados.

\subsection{El problema de la identidad}

Un mismo expediente puede publicarse varias veces: aviso previo, rectificación,
aclaración, formalización o anulación. Cada fuente usa su propio identificador
y sus propias reglas de versionado. Parte central del diseño es una identidad
determinista que permita plegar revisiones sin perder el histórico: en Silver
nada se sobrescribe, todo es un evento con orden y causa.

\begin{table}[ht]
\centering
\small
\begin{tabularx}{\textwidth}{@{}l l X@{}}
\toprule
\cabecera{Fuente} & \cabecera{\texttt{event\_id}} & \cabecera{Regla}\\
\midrule
TED & \texttt{ted:notice:<ND>} & Un número de aviso es un evento; repetirlo no crea filas\\
OpenPLACSP & \texttt{placsp:notice:<atom>@<updated>} & Cada \texttt{updated} válido es una revisión distinta\\
OpenPLACSP & \texttt{placsp:tombstone:<ref>} & El borrado es una fila, no una eliminación\\
\bottomrule
\end{tabularx}
\caption{Identidad determinista por fuente: sin UUIDs ni relojes; reprocesar el
mismo payload produce el mismo \texttt{event\_id}.}
\end{table}

Si aparece contenido distinto para un \texttt{event\_id} ya existente, la
transformación falla e identifica la colisión. Ninguna versión sobrescribe a
la otra de forma silenciosa.

\begin{figure}[ht]
\centering
\begin{tikzpicture}[
  font=\footnotesize,
  evento/.style={draw=ink, line width=0.5pt, fill=white,
    minimum height=1.1cm, minimum width=3.1cm, align=center},
  ultimo/.style={evento, dashed},
  barra/.style={draw=none, fill=paper2,
    minimum height=1.0cm, minimum width=14.3cm, align=center},
  vista/.style={draw=ink, line width=0.5pt, fill=white,
    minimum height=1.1cm, minimum width=7.2cm, align=center},
  flecha/.style={-{Latex[length=2mm]}, line width=0.5pt, color=ink},
  conector/.style={line width=0.5pt, color=muted}
]
\node[evento] (e1) at (1.55,0) {Aviso previo\\ {\scriptsize\color{muted}publicación inicial}};
\node[evento] (e2) at (5.25,0) {Rectificación\\ {\scriptsize\color{muted}corrige importes}};
\node[evento] (e3) at (8.95,0) {Formalización\\ {\scriptsize\color{muted}contrato firmado}};
\node[ultimo] (e4) at (12.65,0) {Anulación\\ {\scriptsize\color{muted}tombstone}};
\draw[flecha] (e1) -- (e2);
\draw[flecha] (e2) -- (e3);
\draw[flecha] (e3) -- (e4);
\node[barra] (silver) at (7.1,-2.0) {Silver · \texttt{procurement\_events}\\ {\scriptsize\color{ink}conserva los 4 eventos con su orden y su causa}};
\draw[conector] (1.55,-0.55) -- (1.55,-1.5);
\draw[conector] (5.25,-0.55) -- (5.25,-1.5);
\draw[conector] (8.95,-0.55) -- (8.95,-1.5);
\draw[conector] (12.65,-0.55) -- (12.65,-1.5);
\node[fill=white, text=muted, inner sep=2pt, font=\scriptsize, align=center] at (7.1,-1.02)
  {una sola identidad determinista\\ \texttt{placsp:procedure:<atom\_id>}};
\draw[flecha] (7.1,-2.5) -- (7.1,-3.1);
\node[vista] (gold) at (7.1,-3.65) {Gold: estado vigente\\ {\scriptsize\color{muted}Anulado · el tombstone no borra el histórico}};
\end{tikzpicture}
\caption{Plegado de revisiones: un expediente, cuatro eventos, una identidad.
Silver conserva todo el histórico; Gold deriva el estado vigente sin destruir
nada.}
\end{figure}

% ============================================================
%  4. ARQUITECTURA
% ============================================================
\section{Arquitectura del pipeline}

El diseño sigue el patrón \emph{medallion}: capas con contrato explícito entre
ellas, de menor a mayor calidad semántica.

\begin{figure}[ht]
\centering
\begin{tikzpicture}[
  font=\footnotesize,
  node distance=0.5cm,
  capa/.style={draw=ink, line width=0.5pt, fill=white,
    minimum height=1.35cm, minimum width=2.2cm, align=center},
  fuente/.style={capa, minimum width=3.0cm},
  p0/.style={dashed, line width=0.5pt, color=muted, -{Latex[length=2mm]}},
  flecha/.style={-{Latex[length=2mm]}, line width=0.5pt, color=ink}
]
\node[fuente] (ted) {TED\\ API JSON\\ {\scriptsize\color{muted}ámbito europeo}};
\node[fuente, below=of ted] (placsp) {OpenPLACSP\\ ZIP · Atom · CODICE\\ {\scriptsize\color{muted}ámbito nacional}};
\node[capa, right=0.9cm of ted.east, yshift=0.125cm] (raw) {Raw\\ JSONL\\ {\scriptsize\color{muted}copia íntegra}};
\node[capa, right=of raw] (bronze) {Bronze\\ Parquet\\ {\scriptsize\color{muted}nada se pierde}};
\node[capa, right=of bronze] (silver) {Silver\\ canónico\\ {\scriptsize\color{muted}eventos}};
\node[capa, below=0.75cm of bronze] (tender) {TenderRecord\\ en memoria\\ {\scriptsize\color{muted}vista legada}};
\node[capa, right=0.5cm of tender] (mid) {Enlace y\\ clasificación\\ {\scriptsize\color{muted}baseline}};
\node[capa, right=0.5cm of mid] (gold) {Gold\\ vista + feed\\ {\scriptsize\color{muted}estado vigente}};
\begin{scope}[on background layer]
  \node[fill=paper2, inner sep=7pt, fit=(raw)(bronze)(silver)(tender)(mid)(gold)] (bg) {};
\end{scope}
\node[anchor=south west, font=\scriptsize\bfseries, color=muted]
  at ([yshift=3pt]bg.north west) {PLATAFORMA · BATCH LOCAL, SIN RED};
\draw[flecha] (ted.east) -- (raw.west);
\draw[flecha] (placsp.east) -- ++(0.5,0) |- (raw.west);
\draw[flecha] (raw) -- (bronze);
\draw[flecha] (bronze) -- (silver);
\draw[flecha] (bronze.south) -- (tender.north);
\draw[flecha] (tender.east) -- (mid.west);
\draw[flecha] (mid.east) -- (gold.west);
\draw[p0] (silver.east) -- ++(0.6,0) -| (gold.north);
\node[anchor=east, font=\scriptsize, color=muted] at (12.9,-0.7) {P0};
\end{tikzpicture}
\caption{Vista inicial v0.1: las fuentes alimentan Raw y Bronze y Silver ya se
construye como canónico, pero el Gold provisional todavía consumía la vista
legada en memoria (TenderRecord) con enlace y clasificación baseline. La
cadena final, con Gold derivado de Silver, se muestra en la Figura 3.}
\end{figure}

\subsection{Del prototipo a la cadena final}

La primera versión completa construía Gold a partir de una vista legada en
memoria que colapsaba revisiones y mezclaba lógica de negocio con ejecución.
Silver pasó a ser el único contrato persistido de entrada: contiene los
eventos canónicos necesarios y permite retirar esa dependencia.

Spark obtiene el estado vigente y aplica el enriquecimiento sobre ese
histórico. DuckDB queda limitado a la consulta analítica de Gold y la
exportación a Tableau conserva un orden determinista. Las fronteras Parquet
permiten reejecutar y auditar cada etapa sin repetir toda la cadena.

\subsection{Decisiones de diseño}

\begin{itemize}
  \item Los registros que no validan van a
        un área de rechazos con su motivo, no desaparecen.
  \item Silver preserva revisiones y anulaciones como
        eventos; el estado vigente es una vista derivada, no el almacenamiento.
  \item El contrato de datos está versionado entre capas; los cambios de esquema
        son cambios de contrato, visibles y revisables.
  \item Cada ejecución emite un manifiesto con
        conteos y checksums de las salidas.
\end{itemize}

\begin{figure}[ht]
\centering
\begin{tikzpicture}[
  font=\footnotesize,
  fuente/.style={draw=ink, line width=0.5pt, fill=white,
    minimum height=1.35cm, align=center},
  barra/.style={draw=ink, line width=0.5pt, fill=white,
    minimum height=1.05cm, minimum width=12.4cm, align=center},
  gris/.style={barra, fill=paper2},
  flecha/.style={-{Latex[length=2mm]}, line width=0.5pt, color=ink}
]
\node[fuente, minimum width=2.2cm] (s1) at (3.0,0) {TED\\ {\scriptsize\color{muted}API JSON}};
\node[fuente, minimum width=2.5cm] (s2) at (5.6,0) {OpenPLACSP\\ {\scriptsize\color{muted}licitaciones}};
\node[fuente, minimum width=2.2cm] (s3) at (8.4,0) {CPV 2008\\ {\scriptsize\color{muted}referencia}};
\node[fuente, minimum width=2.3cm] (s4) at (11.0,0) {DIR3\\ {\scriptsize\color{muted}organismos}};
\node[gris] (raw) at (6.9,-1.85) {Raw inmutable\\ {\scriptsize\color{ink}payload + ventana + timestamp + checksum}};
\node[gris] (bronze) at (6.9,-3.15) {Bronze Parquet\\ {\scriptsize\color{ink}parsing + procedencia + rechazos}};
\node[gris] (silver) at (6.9,-4.45) {Silver Parquet\\ {\scriptsize\color{ink}\texttt{procurement\_events} · histórico canónico}};
\node[gris] (spark) at (6.9,-5.75) {Spark · estado vigente + enriquecimiento\\ {\scriptsize\color{ink}vigencia · CPV · DIR3}};
\node[gris] (gold) at (6.9,-7.05) {Gold Parquet\\ {\scriptsize\color{ink}\texttt{open\_opportunities} · una fila por procedimiento}};
\node[gris] (duck) at (6.9,-8.35) {DuckDB · vistas analíticas\\ {\scriptsize\color{ink}solo lectura sobre Gold · comprador · CPV · mensual}};
\node[barra] (tab) at (6.9,-9.65) {Tableau · dashboard de oportunidades\\ {\scriptsize\color{muted}\texttt{open\_opportunities.csv}}};
\foreach \x in {3.0,5.6,8.4,11.0} \draw[flecha] (\x,-0.675) -- (\x,-1.3);
\draw[flecha] (6.9,-2.42) -- (6.9,-2.6);
\draw[flecha] (6.9,-3.72) -- (6.9,-3.9);
\draw[flecha] (6.9,-5.02) -- (6.9,-5.2);
\draw[flecha] (6.9,-6.32) -- (6.9,-6.5);
\draw[flecha] (6.9,-7.62) -- (6.9,-7.8);
\draw[flecha] (6.9,-8.92) -- (6.9,-9.1);
\end{tikzpicture}
\caption{Arquitectura final implementada: cuatro fuentes oficiales, capas
inmutables con procedencia, histórico canónico en Silver, estado vigente y
enriquecimiento con Spark, Gold en Parquet, vistas analíticas con DuckDB y
consumo en Tableau. La ejecución es batch local; la orquestación programada
queda fuera de esta versión.}
\end{figure}

% ============================================================
%  5. IMPLEMENTACIÓN
% ============================================================
\clearpage
\section{Implementación}

\Needspace{12\baselineskip}
\subsection{Stack}

\begin{table}[H]
\centering
\small
\begin{tabularx}{\textwidth}{@{}l X@{}}
\toprule
\cabecera{Componente} & \cabecera{Elección y motivo}\\
\midrule
Python 3.10+ & Parsers de fuentes, procedencia, Silver y CLI\\
Polars & Procesamiento columnar de la capa Silver\\
Parquet & Formato columnar entre etapas\\
PySpark 4.0.1 & Estado vigente, enriquecimiento y Gold (extra de ejecución)\\
DuckDB & Vistas analíticas solo-lectura sobre Gold\\
Tableau & Export CSV determinista + dashboard de oportunidades\\
uv & Gestión reproducible del entorno y de las ejecuciones\\
CLI (\texttt{licitaciones-pipeline}) & Interfaz única: ingesta, transformación, Gold, analítica y export\\
\bottomrule
\end{tabularx}
\caption{Stack tecnológico final.}
\end{table}

\subsection{Punto de entrada}

La CLI reúne los comandos de la cadena en una única interfaz. La ingesta es la
única etapa que necesita acceso a red; transformación, Gold, analítica y export
se pueden repetir localmente a partir de los artefactos persistidos. El
siguiente listado muestra una ejecución de ejemplo, no el corte de la
ejecución final.

\begin{lstlisting}[caption={Interfaz del pipeline (esquema)}]
# Ingesta explicita (es la unica etapa con red)
licitaciones-pipeline ingest --source ted \
  --start 2026-01-01 --end 2026-06-30
licitaciones-pipeline ingest --source placsp \
  --start 2026-01-01 --end 2026-06-30

# Transformacion sin red y reproducible
licitaciones-pipeline run
licitaciones-pipeline build-gold --as-of 2026-09-17
licitaciones-pipeline build-analytics
licitaciones-pipeline export-tableau
\end{lstlisting}

\subsection{Registro y verificación de procedencia}

Cada fichero Raw viaja con un fichero de procedencia
\texttt{<fichero>.provenance.json}. Este fichero registra la fuente, la ruta,
la ventana temporal, la fecha de recuperación y el checksum. Bronze comprueba
esa evidencia antes de procesar el contenido y detiene la ejecución si falta o
es incoherente.

El corpus histórico era anterior a la introducción de estos ficheros. Para
reutilizarlo sin modificar los originales se añadió una migración de
compatibilidad que comprueba el tamaño y el SHA-256 contra el manifiesto
histórico y crea \emph{hardlinks} en un directorio de trabajo nuevo. El valor
\texttt{run\_at} del manifiesto se registra como instante de observación del
payload, no como hora original de descarga, que no estaba disponible. Esta
limitación queda indicada en \texttt{legacy\_migration\_report.json}.

\Needspace{16\baselineskip}
\begin{lstlisting}[caption={Fichero de procedencia (ejemplo)}]
{
  "version": 1,
  "source": "placsp",
  "raw_path": "placsp/2026-09.zip",
  "sha256": "9f2c...a41b",
  "retrieved_at": "2026-09-02T11:14:03+00:00",
  "partition": "202609",
  "window_start": "2026-09-01",
  "window_end": "2026-09-30"
}
\end{lstlisting}

\Needspace{16\baselineskip}
\subsection{Validación y control de calidad}

La validación empieza en Raw, con la comprobación de la procedencia y el
checksum, y continúa durante el parsing de Bronze, la construcción de la
identidad canónica de Silver y los contratos de Gold. Se comprueban tipos,
claves, fechas, importes, conteos entre capas y unicidad del procedimiento.
Los manifiestos y los informes de ingesta conservan los conteos y el estado de
estas comprobaciones.

Cada registro rechazado se escribe con un valor de
\texttt{rejection\_reason}, por ejemplo \texttt{invalid\_json} o
\texttt{source\_mismatch}. El contenido original permanece en Raw y puede
localizarse mediante su procedencia.

\begin{table}[H]
\centering
\small
\begin{tabularx}{\textwidth}{@{}l X@{}}
\toprule
\cabecera{Salida Bronze} & \cabecera{Contenido}\\
\midrule
\texttt{records.parquet} & Registros aceptados con procedencia completa\\
\texttt{tombstones.parquet} & Cada control de borrado válido, sin deduplicar\\
\texttt{rejections.parquet} & Descartados con motivo (\texttt{rejection\_reason}) y alcance\\
\texttt{ingestion\_report.json} & Conteos: \texttt{parsed = accepted + rejected}, por fuente\\
\bottomrule
\end{tabularx}
\caption{Salidas generadas por la capa Bronze.}
\end{table}

\subsection{Decisión del motor de Silver}

La implementación inicial en Python se utilizó como referencia para comparar
los motores candidatos. Polars y Spark solo se consideraron válidos cuando
produjeron el mismo resultado semántico que esa referencia.

\begin{table}[ht]
\centering
\small
\begin{tabular}{@{}l r r l@{}}
\toprule
\cabecera{Motor} & \cabecera{Tiempo} & \cabecera{Memoria} & \cabecera{Resultado}\\
\midrule
Python & 0,617\,s & 360,96\,MB & Referencia\\
Polars & 0,254\,s & 302,31\,MB & Mismo resultado\\
Spark & 10,010\,s & 1.335,80\,MB & Mismo resultado\\
\bottomrule
\end{tabular}
\caption{Comparativa en el perfil \texttt{small} ($\sim$25.000 filas, mediana
de 3 repeticiones). Informe completo en
\texttt{experiments/silver\_engine\_comparison/}.}
\end{table}

En el perfil \texttt{small}, Polars completó la transformación en 0,254 s
frente a 0,617 s de la referencia, unas 2,4 veces más rápido. Spark mantuvo la
paridad, pero sus 10,010 s no compensaron su coste de arranque en esta etapa.
Silver usa Polars como ruta rápida y conserva la implementación inicial para
los lotes que requieren la ruta conservadora.

\begin{figure}[ht]
\centering
\begin{tikzpicture}[font=\footnotesize, x=1cm, y=1cm]
\node[anchor=east] at (-0.2,0) {python-row};
\node[anchor=east] at (-0.2,-0.75) {polars-native};
\node[anchor=east] at (-0.2,-1.5) {spark-native};
\fill[ink] (0,-0.22) rectangle (2.0,0.22);
\fill[ink] (0,-0.97) rectangle (4.8,-0.53);
\fill[ink] (0,-1.72) rectangle (0.2,-1.28);
\node[anchor=west] at (2.15,0) {1,0\texttimes{}};
\node[anchor=west] at (4.95,-0.75) {2,4\texttimes{}};
\node[anchor=west] at (0.35,-1.5) {0,10\texttimes{}};
\draw[color=line, line width=0.4pt] (0,-2.0) -- (8.0,-2.0);
\foreach \x in {0,2,4,6,8} \draw[color=line, line width=0.4pt] (\x,-2.0) -- (\x,-2.12);
\node[font=\scriptsize\color{muted}] at (4,-2.35) {aceleración frente a la referencia (2\,cm = 1\texttimes{})};
\end{tikzpicture}
\caption{Aceleración frente a la referencia en el perfil \texttt{small}:
Polars $\sim$2,4\texttimes{} con el mismo resultado; Spark local queda en
0,10\texttimes{} en esta escala por su coste de arranque. Tiempos
absolutos en la tabla anterior.}
\end{figure}

% ============================================================
%  6. EVALUACIÓN
% ============================================================
\section{Evaluación}

La evaluación se apoya en un marco reproducible con perfiles de carga de
distinto tamaño y métricas conservadas en JSON. Ninguna cifra de rendimiento
se publica en la memoria sin su ejecución registrada.

\Needspace{8\baselineskip}
\subsection{Criterio de aceptación de resultados}

Una implementación más rápida solo se acepta si produce exactamente los mismos
resultados que la de referencia, de forma consistente en los perfiles medidos.
La medición aceptó combinar Polars con la implementación inicial como respaldo;
Spark no se incorporó a esta etapa del pipeline. Si los resultados difieren,
la divergencia no se descarta: se documenta y se convierte en hallazgo de la
memoria.

Las pruebas de integridad incluyen una colisión deliberada: si dos contenidos
distintos comparten identificador, la transformación debe fallar en lugar de
elegir una versión.

\begin{table}[H]
\centering
\small
\begin{tabularx}{\textwidth}{@{}l r r r X@{}}
\toprule
\cabecera{Perfil} & \cabecera{TED} & \cabecera{PLACSP} & \cabecera{Filas} & \cabecera{Uso}\\
\midrule
\texttt{tiny} & 120 & 120 & $\sim$310 & Pruebas automáticas, milisegundos\\
\texttt{small} & 10.000 & 10.000 & $\sim$25.000 & Comparación local\\
\texttt{medium} & 100.000 & 100.000 & $\sim$250.000 & Escala heurística\\
\texttt{large} & 800.000 & 800.000 & $\sim$2\,M & Exploratorio n=1 (2026-09-16)\\
\texttt{backfill} & 2.000.000 & 2.000.000 & $\sim$5{,}1\,M & Definido, sin resultado conservado\\
\bottomrule
\end{tabularx}
\caption{Perfiles de prueba: conjuntos sintéticos de distinto tamaño.}
\end{table}

\begin{table}[H]
\centering
\small
\begin{tabularx}{\textwidth}{@{}X l l@{}}
\toprule
\cabecera{Criterio} & \cabecera{Se mide con} & \cabecera{Umbral}\\
\midrule
Mismo resultado Bronze\tarrow Silver & Contrato semántico comparado & Igualdad exacta en todos los motores\\
Rendimiento de la etapa & Tiempo y memoria por perfil (JSON) & Mediana de 3 repeticiones\\
Cobertura de ingesta & Conteos del manifiesto frente a fuente & 349.399 entradas PLACSP verificadas\\
Integridad de identidad & Pruebas de colisión y unicidad & Ninguna colisión resuelta en silencio\\
\bottomrule
\end{tabularx}
\caption{Criterios de evaluación del pipeline.}
\end{table}

% ============================================================
%  7. RESULTADOS
% ============================================================
\clearpage
\section{Resultados finales}

\subsection{Ejecución E2E final}

La cadena completa se ejecutó dos veces sobre el corpus de referencia con fecha
de corte 30 de junio de 2026. Ambas ejecuciones partieron de los mismos siete
payloads verificados y recorrieron Bronze, Silver, estado vigente, Gold,
DuckDB y el export para Tableau. Los conteos coincidieron en todas las etapas.

\begin{table}[H]
\centering
\small
\begin{tabularx}{0.72\textwidth}{@{}X r@{}}
\toprule
\cabecera{Etapa} & \cabecera{Filas}\\
\midrule
Bronze aceptadas & 359.304\\
Rechazos Bronze & 0\\
Tombstones Bronze & 1.046\\
Silver canónico & 359.180\\
Estado vigente & 130.118\\
Oportunidades abiertas & 14.763\\
\bottomrule
\end{tabularx}
\caption{Conteos coincidentes en las dos ejecuciones E2E.}
\end{table}

El plegado produjo 130.118 procedimientos vigentes. Gold excluyó 113.957 por
estado cerrado, 521 por borrado y 877 por evidencia insuficiente. Las 14.763
filas restantes cumplen la política de oportunidad abierta para la fecha de
corte. Los 9.905 eventos TED se conservaron en Silver, pero quedaron registrados
como incidencias de estado vigente porque no fue posible resolver una identidad
de procedimiento compatible con el contrato actual.

La reproducibilidad del producto final se comprobó comparando los cinco CSV de
ambas ejecuciones. Cada par tiene el mismo tamaño y el mismo SHA-256, por lo que
su contenido es idéntico byte a byte. Los manifiestos de export difieren solo
en la ruta de la base DuckDB de origen, específica de cada directorio de
ejecución.

\subsection{Productos analíticos y demostración en Tableau}

DuckDB expone vistas de solo lectura sobre Gold y el comando
\texttt{export-tableau} genera cinco CSV con orden determinista. Para demostrar
que estos productos pueden consumirse en una herramienta de BI se seleccionó
uno de ellos, \path{open_opportunities.csv}, y se conectó a una visualización
en Tableau. Este fichero contiene las 14.763 oportunidades, casi 15.000, que
Gold considera abiertas en la fecha de corte del 30 de junio de 2026, con una
fila por procedimiento. Los otros cuatro CSV se conservan como productos
analíticos reproducibles, pero no forman parte de esta demostración.

\begin{table}[H]
\centering
\small
\begin{tabularx}{\textwidth}{@{}l r X@{}}
\toprule
\cabecera{CSV} & \cabecera{Filas} & \cabecera{Granularidad}\\
\midrule
\texttt{open\_opportunities.csv} & 14.763 & Una fila por \texttt{procedure\_id}\\
\texttt{opportunity\_cpv.csv} & 28.730 & Una fila por procedimiento, CPV y posición\\
\texttt{buyer\_summary.csv} & 4.027 & Una fila por identidad analítica del comprador\\
\texttt{cpv\_summary.csv} & 3.751 & Una fila por \texttt{cpv\_code}\\
\bottomrule
\end{tabularx}
\caption{CSV analíticos con datos; la demostración en Tableau utiliza
\texttt{open\_opportunities.csv}.}
\end{table}

% ============================================================
%  8. DISCUSIÓN
% ============================================================
\section{Discusión}

La identidad y el histórico condicionan los resultados. Un procedimiento
puede originar varias publicaciones, por lo que contar filas sin resolver sus
revisiones mezcla eventos y procedimientos. La identidad determinista y las
pruebas de colisión permiten conservar esas revisiones y detectar contenidos
incompatibles sin escoger una versión de forma arbitraria.

La doble ejecución aporta evidencia directa de reproducibilidad. Los conteos
coincidieron desde Bronze hasta Gold y los cinco CSV finales conservaron el
mismo SHA-256. Los ficheros de procedencia, los esquemas versionados y los
manifiestos permiten localizar el origen de esos resultados. El perfil
experimental \texttt{backfill} sigue definido, pero no se presenta como
resultado porque no existe una ejecución conservada.

El benchmark local justificó Polars para Silver: en el perfil
\texttt{small} redujo la transformación de 0,617 s a 0,254 s con el mismo
resultado. Spark fue más lento en esa frontera, pero se utiliza en el estado
vigente y el enriquecimiento, donde se aplican ventanas y joins sobre el
histórico. La elección depende de la operación y del volumen medido en cada
etapa.

La arquitectura final separa dos necesidades. Silver conserva todo el
histórico publicado y Gold ofrece el estado vigente para análisis. Esta
separación evita destruir evidencia al actualizar una vista de negocio.

Los resultados también revelan dos límites semánticos. La identidad de los
9.905 eventos TED no puede trasladarse todavía al grano por procedimiento del
estado vigente. Además, ninguna de las oportunidades abiertas contiene fecha
de publicación. Mantener el campo nulo evita inventar información, pero impide
construir la serie mensual. Ambos problemas están en la normalización anterior
a Gold y no se resuelven filtrando las oportunidades válidas.

El corpus TED usado como referencia está sesgado hacia tecnología por la
consulta de adquisición. Esa limitación afecta a los resultados descriptivos
del corpus, pero no restringe el modelo de datos: la ingesta, las entidades
canónicas y el enriquecimiento CPV admiten contratación de cualquier sector.

% ============================================================
%  9. CONCLUSIONES
% ============================================================
\clearpage
\section{Conclusiones y trabajo futuro}

El proyecto integra fuentes heterogéneas de contratación pública en una cadena
trazable. La decisión principal fue representar cada publicación como un
evento con identidad determinista. Así se conservan las correcciones y
anulaciones y se puede derivar el estado vigente sin perder el histórico.

Raw conserva la evidencia original, Bronze normaliza y registra rechazos,
Silver mantiene el histórico canónico y Gold publica las oportunidades
vigentes. DuckDB y los CSV de Tableau completan la frontera de consumo sin
duplicar la lógica de negocio. Dos ejecuciones completas produjeron los mismos
conteos y cinco CSV idénticos byte a byte. El resultado demuestra la
reproducibilidad de la cadena para el corpus y la configuración evaluados.

El trabajo pendiente se concentra en la semántica de origen: resolver la
identidad de procedimiento de TED y recuperar la fecha de publicación cuando
esté disponible en la fuente. La ausencia de fecha se mantiene como nulo y no
se utiliza para descartar oportunidades abiertas.

\begin{table}[ht]
\centering
\small
\begin{tabularx}{\textwidth}{@{}l >{\raggedright\arraybackslash}X >{\raggedright\arraybackslash}X@{}}
\toprule
\cabecera{Área} & \cabecera{Hoy} & \cabecera{Corrección prevista}\\
\midrule
Corpus & Ventana ene--jun 2026; TED sesgado a tecnología & Ampliar histórico y filtros sectoriales\\
Identidad TED & 9.905 eventos sin identidad de procedimiento resoluble & Revisar la extracción y el puente de identidad\\
Fechas & 14.763 oportunidades sin \texttt{publication\_date} & Recuperar la fecha upstream cuando exista\\
Ejecución & Batch local reproducible y empaquetado con Docker; sin orquestación & Programación con Airflow\\
Dashboard & Conectado solo a \texttt{open\_opportunities.csv} & Añadir otros conjuntos solo si el caso de uso lo exige\\
\bottomrule
\end{tabularx}
\caption{Limitaciones vigentes y líneas futuras: ingesta periódica,
enriquecimiento con nuevas fuentes (p.\,ej.\ contratos menores) y semántica
para ranking y búsqueda.}
\end{table}

% ============================================================
%  BIBLIOGRAFÍA
% ============================================================
\clearpage
\begin{thebibliography}{99}
\addcontentsline{toc}{section}{Bibliografía}

\bibitem{ted} Publications Office of the European Union.
\emph{TED eSender: API documentation}. Oficina de Publicaciones de la UE.

\bibitem{placsp} Ministerio de Economía, Comercio y Empresa.
\emph{Open Data Portal de la Plataforma de Contratación del Sector Público}.

\bibitem{dir3} Ministerio de Políticas Territoriales.
\emph{DIR3: Directorio Común de unidades orgánicas y oficinas}.

\bibitem{cpv} Reglamento (UE) 2015/2424, que modifica el Reglamento (CE) 2195/2007
sobre el Vocabulario Común de Contratos Públicos (CPV).

\bibitem{polars} Polars. \emph{User guide}. \url{https://pola.rs}

\bibitem{spark} Apache Software Foundation. \emph{Apache Spark documentation}.
\url{https://spark.apache.org/docs/latest/}

\bibitem{duckdb} DuckDB Foundation. \emph{DuckDB documentation}.
\url{https://duckdb.org/docs/stable/}

\bibitem{tableau} Salesforce. \emph{Tableau documentation}.
\url{https://help.tableau.com/}

\bibitem{repositorio} Campoy Rojo, T. (2026).
\emph{tfm-licitaciones: pipeline DataOps reproducible para contratación
pública}. Repositorio de código en GitHub.
\url{https://github.com/tomcrojo/tfm-licitaciones}. Consulta: 17 de septiembre
de 2026.

\bibitem{kleppmann} Kleppmann, M. (2017). \emph{Designing Data-Intensive
Applications}. O'Reilly.

\end{thebibliography}

\end{document}
